% CV completo — Español. Base: trayectoria completa + aportes de IA y founder.
\documentclass[11pt,a4paper]{article}
\usepackage[T1]{fontenc}
\usepackage[utf8]{inputenc}
\usepackage[margin=1.4cm]{geometry}
\usepackage{helvet}
\renewcommand{\familydefault}{\sfdefault}
\usepackage{enumitem}
\usepackage{titlesec}
\usepackage{xcolor}
\usepackage[hidelinks]{hyperref}

\definecolor{ink}{HTML}{111111}
\definecolor{muted}{HTML}{555555}
\color{ink}
\pagestyle{empty}
\setlist[itemize]{leftmargin=1.1em,itemsep=1pt,topsep=2pt}

\titleformat{\section}{\large\bfseries\scshape}{}{0em}{}[{\color{ink}\titlerule[1.2pt]}]
\titlespacing{\section}{0pt}{1.0em}{0.45em}

\newcommand{\job}[3]{\vspace{3pt}\noindent{\bfseries #1} \textcolor{muted}{— #2}\hfill\textcolor{muted}{\small #3}\par}
\newcommand{\heads}[1]{\textbf{#1:} }

\begin{document}

{\Huge\bfseries Danrley Pereira}\\[2pt]
{\large\color{muted} Fundador e Ingeniero de Software Full-Stack — IA, Backend y Producto}\\[4pt]
{\small
danrley.pereira@dwcorp.com.br \textbullet{} +55 61 9 9423-4712 \textbullet{} Brasília, Brasil \textbullet{}
\href{https://www.linkedin.com/in/danrleypereira}{LinkedIn} \textbullet{}
\href{https://github.com/danrleypereira}{GitHub} \textbullet{}
\href{https://danrleypereira.com}{danrleypereira.com}
}

\section{Resumen}
Ingeniero de software full-stack desde 2017 y fundador de DW Corp, construyendo \textbf{productos habilitados por IA de extremo a extremo} — agentes LLM, pipelines RAG y machine learning — sobre una base sólida de backend en Python, C\#/.NET y Node.js. He llevado funciones de IA a clientes reales, construido \textbf{bases de producto reutilizables} (autenticación, pagos, CI/CD, infraestructura en la nube y on-premise), contratado y mentorizado ingenieros, y lidero la primera comunidad de \textbf{Software Craftsmanship} de Brasília. Soy responsable de todo el camino, desde los datos y la recuperación hasta los servicios de backend y las interfaces web que los entregan, y tengo amplia experiencia en software operativo y de logística.

\section{Habilidades Principales}
\begin{itemize}
\item \heads{IA / ML} APIs de LLM (OpenAI, Anthropic, Azure OpenAI), RAG y embeddings, FAISS, LangChain / LangGraph, flujos con agentes, prompt engineering, TensorFlow / Keras, aprendizaje por refuerzo.
\item \heads{Backend} Python (FastAPI, Flask, Django, Celery), C\#/.NET y ASP.NET, Node.js, Rust, APIs RESTful, arquitectura orientada a eventos (RabbitMQ), microservicios.
\item \heads{Frontend} React, Next.js, Vue / Nuxt, Angular, TypeScript.
\item \heads{Datos} PostgreSQL, SQLAlchemy, MySQL, MongoDB, Redis, Snowflake, ETL / pipelines de datos, Pandas.
\item \heads{Nube y DevOps} Docker, Kubernetes, AWS, GCP, Azure, CI/CD (GitHub y GitLab), Terraform, Grafana, on-premise / Linux bare-metal.
\item \heads{Pruebas} Pytest, Playwright, Cypress, Cucumber, Behave, Allure, TDD / BDD.
\end{itemize}

\section{Experiencia}

\job{Fundador e Ingeniero de IA}{DW Corp}{sep 2022 -- Actualidad}
\begin{itemize}
\item Construí \textbf{agentes LangChain / LangGraph} — un asistente de presupuesto por WhatsApp y un bus multi-agente orientado a eventos — tras una interfaz común multi-proveedor de LLM.
\item Construí \textbf{Curupira}, una plataforma de autenticación Rust/OIDC reutilizable (multi-tenancy, SSO, MFA), además de integraciones de pago (Stripe, Adyen) y un stack bare-metal + nube híbrida con Terraform y Grafana.
\item Entregué productos para clientes: un gestor de pedidos de restaurantes (Django, Celery, resiliencia con circuit-breaker hacia un data lake) y una app de field-service en .NET MAUI con procedimientos, vistas y triggers de base de datos.
\item Contraté y mentoricé a más de 10 ingenieros, y lidero la primera comunidad de Software Craftsmanship de Brasília (katas de TDD, mesas redondas).
\end{itemize}

\job{Ingeniero de Software Senior}{Epicor}{nov 2023 -- mar 2025}
\begin{itemize}
\item Integré \textbf{Azure OpenAI} y un \textbf{pipeline RAG} sobre datos propietarios; usé Azure Document Intelligence para extracción estructurada.
\item Construí microservicios \textbf{.NET Aspire} con procesamiento orientado a eventos (\textbf{RabbitMQ}), OAuth2/SAML personalizado, desplegados en Kubernetes.
\item Lideré la refactorización de un portal Angular a microfrontends, con sólida cobertura xUnit y end-to-end, e impulsé la planificación de sprints y la documentación.
\end{itemize}

\job{Investigador}{Ibict / MCTI (gobierno federal)}{mar 2025 -- Actualidad}
\begin{itemize}
\item Construí pipelines de datos en \textbf{Snowflake} sobre datos institucionales, avanzando hacia el entrenamiento de modelos con datos propietarios.
\item Co-orienté una tesis de grado que diseñó un \textbf{sistema RAG} para el poder judicial brasileño (calificada 9,0/10); audité una plataforma de más de 600 mil usuarios frente a ISO/IEC 25000 y LGPD.
\end{itemize}

\job{Consultor Full-Stack Senior}{SubSurface Dynamics}{ene 2024 -- jun 2024}
\begin{itemize}
\item Refactoricé una app MERN a un proyecto \textbf{TypeScript} moderno sobre \textbf{Fastify, Socket.io, Redis} y Typegoose, sirviendo a miles de usuarios en tiempo real.
\item Ajusté Redis y las series temporales de MongoDB Atlas por costo y rendimiento; aseguré la plataforma con Azure AD (ADAL) + Auth0 (OAuth2/SAML).
\item Construí gráficos 2D/3D y una visualización de pozos con \textbf{Three.js}, además de modo oscuro y un sistema de notificaciones.
\end{itemize}

\job{Ingeniero de Datos y DevOps}{Dado Global}{dic 2021 -- abr 2023}
\begin{itemize}
\item Diseñé un sistema \textbf{ETL} con responsabilidades por capas usando \textbf{SQLAlchemy} y cobertura rigurosa de pytest; automaticé la extracción con Selenium y Behave.
\item Llevé la entrega a \textbf{Continuous Delivery} con GitHub Actions (semanal $\rightarrow$ diaria), mejorando el SLA del 70\% al 95\%.
\end{itemize}

\job{Ingeniero de Software}{ClearSale}{nov 2021 -- mar 2022}
\begin{itemize}
\item Construí microservicios antifraude de alto rendimiento en \textbf{C\#/.NET}; diseñé vistas, triggers y procedimientos almacenados y optimicé consultas con Dapper.
\item Definí procesos de ingeniería (UML/BPMN), inicié pruebas unitarias y mentoricé al equipo.
\end{itemize}

\job{Líder Técnico e Investigador}{LabTech, UDF}{ene 2023 -- dic 2023}
\begin{itemize}
\item Lideré una plataforma multinivel de gestión de eventos (React, Python, Spring Boot, RabbitMQ, Kubernetes) y mentoricé a ingenieros junior.
\end{itemize}

\job{Ingeniero de Software}{Layers Education}{mar 2022 -- sep 2022}
\begin{itemize}
\item Mantuve sistemas backend y microfrontend; reduje los tiempos de consulta de MongoDB de 1,3s a 0,6s bajo carga, construí un data lake en BigQuery para ETL y reforcé la seguridad según OWASP.
\end{itemize}

\job{Ingeniero Full-Stack}{Compuletra}{jul 2020 -- ago 2021}
\begin{itemize}
\item Construí software de logística de \textbf{transporte, entrega y patio} con sistemas de flota e \textbf{inspección de vehículos}, y una plataforma ERP de escritorio (Electron) sobre PostgreSQL.
\item Migré una app legacy de seguimiento vehicular de AngularJS a React (on-premise), y construí un generador de páginas comerciales con checkout compatible con PCI gestionado desde la base de datos.
\end{itemize}

\job{Desarrollador Full-Stack y BPM}{ConsultMídia / Banco de Brasília}{abr 2020 -- ene 2021}
\begin{itemize}
\item Construí interfaces AngularJS/Spring Boot y APIs REST sobre BPM \textbf{Camunda} para modernizar los flujos del banco; integré Alfresco (gestión documental) y automaticé pruebas de rendimiento (JMeter, Selenium).
\end{itemize}

\job{Desarrollador Full-Stack y Móvil}{Transoft}{feb 2021 -- oct 2021}
\begin{itemize}
\item Construí y renové aplicaciones para gestionar sensores y cámaras de flotas de autobuses, y una app móvil de gestión de personal integrada con Oracle (Angular, Ionic, Laravel).
\end{itemize}

\job{Desarrollador de Plataforma}{Ministerio de Salud / USP}{jun 2019 -- jul 2020}
\begin{itemize}
\item Ayudé al Ministerio de Salud a gestionar la \textbf{distribución de medicamentos en más de 5.000 municipios}; construí el frontend e integré APIs REST y microservicios para exponer datos, reducir el fraude y optimizar la integración de servicios.
\end{itemize}

\section{Experiencia Anterior}
\begin{itemize}
\item \textbf{Taurus Software} — aplicaciones de gestión de energía IoT, móviles híbridas (Ionic) y de captación de leads. \emph{2018--2020}
\item \textbf{Gama Cidadão} — Gerente Web y de Marketing Digital; portales WordPress/Vue, landing pages y SEO. \emph{2017--2018}
\item \textbf{UMESB} — Desarrollador por contrato (sistema de registro de estudiantes, Polymer/Node.js). \emph{2017}
\item \textbf{IT Universe} — Instructor de diseño web/gráfico y fundamentos de computación. \emph{2017}
\end{itemize}

\section{Educación}
\begin{itemize}
\item Grado en Ingeniería de Software — Universidad de Brasília
\item Grado en Ciencias de la Computación — Centro Universitario del Distrito Federal
\item MBA, Cadena de Suministro y Logística Integrada — Southern Cross University
\item MBA, Administración de Bases de Datos — Anhanguera
\end{itemize}

\section{Certificaciones y Charlas}
\begin{itemize}
\item Ponente en el escenario principal de \textbf{Campus Party Brasil} — Clean Architecture, TDD, IA, SDET y Software Craftsmanship.
\item Business Intelligence con Power BI \textbullet{} Scrum Fundamentals (SFC) \textbullet{} Kanban Foundation \textbullet{} Producción Oral B2 (Inglés) — UnB.
\end{itemize}

\section{Idiomas}
Portugués (nativo) \textbullet{} Inglés (C2) \textbullet{} Español (intermedio)

\end{document}
